\section{Proofs for the planar walk}\label{app:planar}

This appendix has three parts. We first prove the bound on $o^{(4)}_N$ in
Proposition~\ref{prop:decomp} and Theorem~\ref{thm:planar-first}. Then we prove the local limit of
Theorem~\ref{thm:fourdir}, and finally parts (b) and (c) of Theorem~\ref{thm:pureturns}. We identify $\R^2$ with $\mathbb C$, so $e_z=\mathrm i^z$ and
$Y_N=\sum_{t=0}^m\ell_t\mathrm i^{Z_t}$ in the notation of Lemma~\ref{lem:threelevel}. Let
$\mathrm{bin}(m)=\binom{N-1}m(\beta h)^mp_0^{N-1-m}$ be the law of $m$, with mean
$\bar m=(N-1)\beta h$. We write $\Prob_z$ and $\E_z$ for the skeleton walk of
Lemma~\ref{lem:threelevel}(c) started at $z$, $\Prob$ for its uniform start, and $\xi$ for one of
its increments. Put $\varkappa=s/\beta$, so $\varkappa<1$ when $r>0$. For a skeleton, let $a,b,c,d$
be the numbers of runs going east, west, north and south. So $a+b+c+d=m+1$, and the path uses all four
directions exactly when $a,b,c,d\ge1$. Throughout, $N=2\nu$, and $[\,\cdot\,]$ is the indicator of
the statement inside.

For $k\ge1$ let $\chi_k(X)=\binom{X-1}{k-1}[X\ge1]$, and let $\chi_0(X)=[X=0]$. Then $\chi_k(X)$
counts the compositions of $X$ into $k$ positive parts.

\begin{lemma}[return counts]\label{lem:returncount}
Let $Z$ be a skeleton with $m$ switches and counts $a,b,c,d$. The number of compositions $\ell$ of
$N$ into $m+1$ parts with $\sum_t\ell_t\mathrm i^{Z_t}=0$ is
\[
 \Omega_N(a,b,c,d)=\sum_{X=0}^{\nu}\chi_a(X)\chi_b(X)
 \chi_c(\nu-X)\chi_d(\nu-X).
\]
If $a,b,c,d\ge1$, then $m\ge3$ and $\Omega_N(a,b,c,d)\le\frac12\binom N{m-2}$.
\end{lemma}

\begin{proof}
The sum vanishes exactly when the runs going east and west have the same total length $X$, and
those going north and south have total length $\nu-X$. The four compositions of these totals into runs
are independent, which gives the formula. Now let $a,b,c,d\ge1$. If $m=3$, there is one run
in each direction and $\Omega_N=\nu-1<\frac12\binom N1$. Let $m\ge4$. Pick $t_z$ with
$Z_{t_z}=z$ for each $z$, and fix the other $m-3$ run lengths. If they add up to $N-j$, there are
$\binom{N-j-1}{m-4}$ ways to choose them. The four free lengths add up to $j$, and $Y_N=0$ fixes
$\ell_{t_0}-\ell_{t_2}$ and $\ell_{t_1}-\ell_{t_3}$. So $\ell_{t_2}$ determines the other three, and
$\ell_{t_0}+\ell_{t_2}$ lies in $[2,j-2]$ with a fixed parity. So there are at most $j/2$ solutions,
and
\[
 \Omega_N(a,b,c,d)\le\sum_{j=0}^{N-1}\frac j2\binom{N-1-j}{m-4}=\frac12\binom N{m-2},
\]
where the identity counts the subsets of $\{0,\dots,N-1\}$ of size $m-2$ by their second smallest
element $j$.
\end{proof}

By Lemma~\ref{lem:threelevel}(b) and (d), $\varpi_N=\Omega_N(a,b,c,d)/\binom{N-1}m$ is the probability that
$Y_N=0$ given $m$ and $Z$, so
\begin{equation}\label{eq:app-planar-o4}
 o^{(4)}_N=\sum_{m\ge3}\mathrm{bin}(m)\,\E\bigl[\varpi_N\mathbf 1\{a,b,c,d\ge1\}\bigm|m\bigr].
\end{equation}

\begin{lemma}[skeleton returns]\label{lem:skelreturn}
Let $W_m=\sum_{t=0}^m\mathrm i^{Z_t}=(a-b)+\mathrm i(c-d)$.
\textup{(a)} $\E_z\lvert W_m\rvert^2\le m+1$ for every start $z$.
\textup{(b)} If $s>0$ and $m$ is odd, then $\Prob(W_m=0)\ge\varkappa/(36(m+1))$.
\textup{(c)} If $s=0$ and $m\equiv3\pmod4$, then $\Prob(W_m=0)\ge1/(m+1)$.
\end{lemma}

\begin{proof}
(a) Since $\E\,\mathrm i^\xi=\frac r\beta(\mathrm i+\mathrm i^{-1})+\frac s\beta\mathrm i^2=-\varkappa$,
we get $\E_z\lvert W_m\rvert^2=\sum_{t,t'}\E_z\mathrm i^{Z_t-Z_{t'}}
=m+1+2\sum_{k=1}^m(m+1-k)(-\varkappa)^k$. The last sum alternates in sign, starts negative and
has nonincreasing absolute terms, so it is at most $0$.

(b) Write $m=2m_1+1$ and split $W_m=W'+W''$, where $W'$ sums over $t\le m_1$ and $W''$ over
$t>m_1$. The law of $\xi$ is symmetric and the uniform law is stationary, so
$(Z_{m_1},\dots,Z_0)$ is again the walk from a uniform start, and
$\Prob(W'=w,\,Z_{m_1}=z)=\frac14\Prob_z(W_{m_1}=w)$. With probability $\varkappa$, independently of
$Z_0,\dots,Z_{m_1}$, the next switch is a reversal. Then given $Z_{m_1}=z$, the second half is the
walk from $z+2$, and $\Prob_{z+2}(W_{m_1}=-w)=\Prob_z(W_{m_1}=w)$. So
$\Prob(W_m=0)\ge\frac\varkappa4\sum_z\sum_w\Prob_z(W_{m_1}=w)^2$. Let $R^2=m+1$. By (a) and
Chebyshev's inequality, $\Prob_z(\lvert W_{m_1}\rvert\le R)\ge\frac12$. At most $9R^2$ Gaussian
integers have modulus at most $R$, so by Cauchy--Schwarz the inner sum is at least $1/(36R^2)$.
Summing over $z$ gives (b).

(c) Now every switch is a turn, so $\mathrm i^{Z_t}=\mathrm i^{Z_{t-1}}\cdot\mathrm i\epsilon_t$ with
independent fair signs $\epsilon_t$. Hence $\mathrm i^{Z_t}=\mathrm i^{Z_0}\mathrm i^t\sigma_t$, where
$\sigma_0=1$ and $\sigma_t=\epsilon_1\cdots\epsilon_t$ are independent fair signs for $t\ge1$. Write
$m=4k-1$. Then $\mathrm i^{-Z_0}W_m$ has real part $\sum_{t\text{ even}}(-1)^{t/2}\sigma_t$, which is
$1$ plus a sum of $2k-1$ fair signs, and an independent imaginary part that is a sum of $2k$ fair
signs. So
\[
 \Prob(W_m=0)=\binom{2k-1}k\binom{2k}k2^{1-4k}=\prod_{j=1}^k\frac{(2j-1)^2}{(2j)^2}
 \ge\frac14\prod_{j=2}^k\frac{2j-2}{2j}=\frac1{4k}=\frac1{m+1}.
\]
\end{proof}

\begin{lemma}[balanced skeletons]\label{lem:balanced}
If $a=b\ge1$, $c=d\ge1$ and $m=2(a+c)-1\le N-1$, then
$\Omega_N(a,b,c,d)/\binom{N-1}m\ge\sqrt{m+1}/(2N^2)$.
\end{lemma}

\begin{proof}
Put $k=a+c=(m+1)/2$, so $2\le k\le\nu$. Only $1\le X\le\nu-1$ contribute to $\Omega_N$, so by
Cauchy--Schwarz and Vandermonde's identity
\[
 \Omega_N=\sum_{X=1}^{\nu-1}\binom{X-1}{a-1}^2\binom{\nu-X-1}{c-1}^2
 \ge\frac1\nu\binom{\nu-1}{k-1}^2=\frac k{\nu^2}\binom{\nu-1}{k-1}\binom\nu k .
\]
Also $\binom{\nu-1}{k-1}\binom\nu k/\binom{2\nu-1}{2k-1}=\Prob(H=k-1)$, where $H$ is the number of
red balls among $2k-1$ balls drawn without replacement from $\nu-1$ red and $\nu$ blue ones. If
$k=\nu$, then $H=k-1$. Otherwise the ratio
$\Prob(H=j+1)/\Prob(H=j)=(\nu-1-j)(2k-1-j)/((j+1)(\nu-2k+2+j))$ decreases in $j$, is
$\frac{\nu-k}{\nu-k+1}<1$ at $j=k-1$ and exceeds $1$ at $j=k-2$. So $k-1$ is a mode of $H$. Since
$\Var H\le(2k-1)/4$, Chebyshev's inequality puts mass at least $\frac34$ on the at most
$3\sqrt{2k}$ integers within $\sqrt{2k-1}$ of $\E H$, so $\Prob(H=k-1)\ge1/(4\sqrt{2k})$. As
$\binom{N-1}m=\binom{2\nu-1}{2k-1}$, we get $\Omega_N/\binom{N-1}m\ge k/(4\nu^2\sqrt{2k})=\sqrt{m+1}/(2N^2)$.
\end{proof}

\begin{proof}[Proof of the bound on $o^{(4)}_N$ in Proposition~\ref{prop:decomp}]
Note that $(\beta h)^mp_0^{N-1-m}=p_0^{N-1}w^m$. By \eqref{eq:app-planar-o4} and
Lemma~\ref{lem:returncount},
\[
 o^{(4)}_N\le\frac{p_0^{N-1}}2\sum_{m=3}^{N-1}\binom N{m-2}w^m
 \le\frac{p_0^{N-1}w^2}2\bigl((1+w)^N-1\bigr)=\frac{w^2}{2p_0}\bigl(1-p_0^N\bigr),
\]
since $1+w=1/p_0$. Dividing by $c_N$ gives the bound on $E_N$.
\end{proof}

\begin{proof}[Proof of Theorem~\ref{thm:planar-first}]
We first prove three bounds. Since $I_1\le I_0\le e^z$ for $z\ge0$, we have $\psi(z)\le2ze^z$, and
\eqref{eq:planar-bessel} gives $2S_N(u)\le8ue^{Nu}$. With Proposition~\ref{prop:decomp} and
$(1+w)^N\le e^{Nw}$,
\begin{equation}\label{eq:app-planar-upper}
 \frac{o_N}{c_N}\le8ue^{Nu}+2w^2e^{Nw}.
\end{equation}
Since $o_N\ge o^{(0)}_N$, \eqref{eq:planar-bessel} with $z=Nu$ gives
\begin{equation}\label{eq:app-planar-axes}
 \frac{o_N}{c_N}\ge2S_N(u)\ge\frac4N\Bigl(1-\frac{z(z+2)}{2N}\Bigr)\psi(z).
\end{equation}
Third, we show that there are $c'>0$ and $T_0$, depending only on $r$ and $s$, with
\begin{equation}\label{eq:app-planar-four}
 o^{(4)}_N\ge\frac{c'}{N^2\sqrt T}\qquad\text{whenever }T=Nh\ge T_0\text{ and }\beta h\le\tfrac12 .
\end{equation}
Let $\mathcal M$ be the set of odd $m$ if $s>0$, and the set of $m\equiv3\pmod4$ if $s=0$. Let
$m\in\mathcal M$. The event $\{a=b\ge1,\ c=d\ge1\}$ is $\{W_m=0\}$ minus $\{a=b=0\}$ and
$\{c=d=0\}$. Each of these needs every switch to be a reversal, so it has probability at most
$\frac12\varkappa^m$, and it is empty if $s=0$. So Lemma~\ref{lem:skelreturn} gives $m_2$ and $c_2>0$
with $\Prob(a=b\ge1,\ c=d\ge1\mid m)\ge c_2/(m+1)$ for $m\in\mathcal M$ with $m\ge m_2$ (take
$c_2=\varkappa/72$ if $s>0$ and $c_2=1$ if $s=0$). On this event Lemma~\ref{lem:balanced} gives
$\varpi_N\ge\sqrt{m+1}/(2N^2)$. Let $J=[\bar m/2,2\bar m]$. By \eqref{eq:app-planar-o4}, if
$\bar m\ge2m_2$ then $o^{(4)}_N\ge\sum_{m\in\mathcal M\cap J}\mathrm{bin}(m)\,c_2/(2N^2\sqrt{m+1})
\ge c_2\Prob(m\in\mathcal M\cap J)/(2N^2\sqrt{2\bar m+1})$. The Chernoff bounds for the
binomial law are $\Prob(m\le\bar m/2)\le e^{-\bar m/8}$ and $\Prob(m\ge2\bar m)\le e^{-\bar m/3}$,
so $\Prob(m\notin J)\le2e^{-\bar m/8}$. Put $p=\beta h\le\frac12$. Then
$\Prob(m\text{ odd})=\frac12(1-(1-2p)^{N-1})\ge\frac12(1-e^{-2\bar m})$. For $s=0$ the roots of unity
filter gives
\[
 \Prob(m\equiv3\bmod4)=\frac14\sum_{k=0}^3\mathrm i^{-3k}\bigl(1-p+p\,\mathrm i^k\bigr)^{N-1}
 \ge\frac14-\frac34e^{-\bar m/2},
\]
since $\lvert1-p+p\,\mathrm i^{\pm1}\rvert^2=1-2p(1-p)\le e^{-p}$ and $\lvert1-2p\rvert\le e^{-2p}$.
So $\Prob(m\in\mathcal M\cap J)\ge\frac18$ when $\bar m$ is large. As $\bar m\ge\beta T/2$, and
$2\bar m+1\le3\beta T$ when $\beta T\ge1$, this proves \eqref{eq:app-planar-four} with
$c'=c_2/(16\sqrt{3\beta})$.

(a) Let $h=\tau L/N$ and $T=\tau L$. Then $u,w=O(L/N)$, $Nu=s\tau L+O(L^2/N)$ and
$Nw=\beta\tau L+O(L^2/N)$. So \eqref{eq:app-planar-upper} gives
$o_N/c_N=O(LN^{s\tau-1}+L^2N^{\beta\tau-2})$. If $\tau<\tau^*$, then $s\tau<1$ and $\beta\tau<2$,
and this is at most $N^{-\gamma}$ for large $N$ whenever $0<\gamma<\min(1-s\tau,2-\beta\tau)$. Now
let $\tau>\tau^*$, so $s\tau>1$ or $\beta\tau>2$. If $s\tau>1$, then $z=Nu\to\infty$ and
$z^2/N\to0$. As $I_0(z)+I_1(z)\sim2e^z/\sqrt{2\pi z}$~\cite[Eq.~10.40.1]{dlmf}, we have
$\psi(z)\sim\sqrt{2z/\pi}\,e^z$, and \eqref{eq:app-planar-axes} gives $o_N/c_N\ge N^{s\tau-1}$ for
large $N$. If $\beta\tau>2$, then \eqref{eq:app-planar-four} and
$c_N^{-1}=4p_0^{-(N-1)}\ge4e^{\beta T-\beta h}$ give
$o_N/c_N\ge4c'e^{-\beta h}N^{\beta\tau-2}(\tau L)^{-1/2}$, which is at least $N^{(\beta\tau-2)/2}$
for large $N$.

(b) Let $\tau<\tau^*<\tau'$. By (a), for large $N$ the ratio $o_N/c_N$ is below $1$ at
$h=\tau L/N$ and above $1$ at $h=\tau'L/N<1/\beta$. It increases in $h$ by
Proposition~\ref{prop:decomp}, so $\tau L<T^*_N<\tau'L$.

(c) Put $\tau_N=T^*_N/L$, so $\tau_N\to\tau^*$ by (b), and $w^*=\beta h^*_N/(1-\beta h^*_N)$. As in
(a), $u^*_N,w^*=O(L/N)$, $Nu^*_N=s\tau_NL+o(1)$ and $Nw^*=\beta\tau_NL+o(1)$. At the crossing
Proposition~\ref{prop:decomp} gives $2S_N(u^*_N)+E_N=1$. If $s>2r$, then $\tau^*=1/s$ and
$\beta\tau^*<2$. So $E_N\le2w^{*2}e^{Nw^*}=N^{\beta/s-2+o(1)}\to0$, and $2S_N(u^*_N)\to1$. The two
axes carry the same mass by symmetry. If $s<2r$, then $\tau^*=2/\beta$ and $s\tau^*<1$. So
$2S_N(u^*_N)\le8u^*_Ne^{Nu^*_N}=N^{2s/\beta-1+o(1)}$, which is at most $N^{-\gamma}$ for large $N$
whenever $0<\gamma<1-2s/\beta$.
\end{proof}

\subsection{The four-direction local limit}

\emph{Idea of the proof.} Given the skeleton, the run lengths are close to the gaps between $m$
uniform points in $[0,N]$.
So the total lengths in the four directions are close to $N$ times a Dirichlet vector
$(y_0,y_1,y_2,y_3)$ with parameters $(a,c,b,d)$. Then $N^2\varpi_N$ should be close to twice the
density of $(y_0-y_2,y_1-y_3)$ at $0$, where the factor $2$ is the parity of $Y_N$. This is the
quantity $g(a,b,c,d)$ below. Divided by $m$, it is close to a Gaussian function of $W_m/\sqrt m$,
and $W_m$ satisfies a central limit theorem. Steps 1 and 2 below remove the skeletons with too few
or too many switches and those that use some direction rarely. Step 3 compares $N^2\varpi_N$ with
$g$, Step 4 averages $g$ over the skeleton, and Step 5 puts the steps together.

\begin{proof}[Proof of Theorem~\ref{thm:fourdir}]
We use the notation of the start of this appendix. Let $A=a+b$ and $B=c+d$, so $A+B=m+1$, and put $\Gamma_4=(a-1)!\,(b-1)!\,(c-1)!\,(d-1)!$ and
\[
 g(a,b,c,d)=\frac{m(m-1)}2\cdot\frac{\binom{A-2}{a-1}}{2^{A-2}}\cdot\frac{\binom{B-2}{c-1}}{2^{B-2}}
 =\frac{m!\,(A-2)!\,(B-2)!}{2^{m-2}\,(m-2)!\,\Gamma_4}.
\]
Let $J_N=[\bar m/2,2\bar m]$, $\mathcal G_m=\{a,b,c,d\ge m/8\}$ and $m_0=20$. Note that
$\bar m=\beta T(1-1/N)\to\infty$ and $\bar m\le\beta N^{1/8}$. Below, $C$ is a constant that depends
only on $r$ and $s$ and may change from line to line.

\emph{Step 1 (few or many switches).} If $a,b,c,d\ge1$ and $m\le N/5$, then
Lemma~\ref{lem:returncount} and
$\binom N{m-2}/\binom{N-1}m=Nm(m-1)/((N-m+2)(N-m+1)(N-m))$ give $\varpi_N\le m^2/N^2$. If $m>N/5$,
then $\varpi_N\le1<25m^2/N^2$. So $N^2\varpi_N\le25m^2$ whenever the path uses all four directions.
As in the proof of Theorem~\ref{thm:planar-first}, the Chernoff bounds give
$\Prob(m\notin J_N)\le2e^{-\bar m/8}$. Also
$\E m^4=\sum_{j=1}^4S(4,j)(N-1)_j(\beta h)^j\le15\bar m^4$, where $S(4,j)=1,7,6,1$ are Stirling
numbers of the second kind, $(N-1)_j$ is a falling factorial, and $\bar m\ge1$. By the
Cauchy--Schwarz inequality, the terms of \eqref{eq:app-planar-o4} with $m\notin J_N$ contribute at
most $25\,\E[m^2\mathbf 1\{m\notin J_N\}]\le25\sqrt{30}\,\bar m^2e^{-\bar m/16}=o(1)$ to
$N^2o^{(4)}_N$.

\emph{Step 2 (unbalanced skeletons).} For $k=1,2,3$ let $V_k=\sum_{t=0}^m\mathrm i^{kZ_t}$. The
number of $t$ with $Z_t=z$ is $\frac{m+1}4+\frac14\sum_{k=1}^3\mathrm i^{-kz}V_k$. Let
$\vartheta_k=\E\,\mathrm i^{k\xi}$, so $\vartheta_1=\vartheta_3=-\varkappa$ and
$\vartheta_2=2\varkappa-1$. Then
$\E[\mathrm i^{kZ_t}\mid Z_0,\dots,Z_{t-1}]=\vartheta_k\mathrm i^{kZ_{t-1}}$, so
$M^{(k)}_m=\sum_{t=1}^m(\mathrm i^{kZ_t}-\vartheta_k\mathrm i^{kZ_{t-1}})$ is a martingale with
increments of modulus at most $2$, and
\[
 (1-\vartheta_k)V_k=M^{(k)}_m+\mathrm i^{kZ_0}-\vartheta_k\mathrm i^{kZ_m}.
\]
Here $1-\vartheta_k\ge\min(1,4r/\beta)>0$. Burkholder's inequality \cite[Theorem~2.10]{hallheyde1980},
applied to the real and imaginary parts, bounds $\E\lvert M^{(k)}_m\rvert^4$ by a constant times
the second moment of the sum of the $m$ squared increments. So
$\E\lvert M^{(k)}_m\rvert^4\le Cm^2$. So
$\E\lvert V_k\rvert^4\le Cm^2$ and $\E(a-\frac{m+1}4)^4\le Cm^2$, and the same holds for $b,c,d$.
Since $\frac{m+1}4-\frac m8\ge\frac m8$, Markov's inequality gives
$\Prob(\mathcal G_m^c\mid m)\le C/m^2$. Every skeleton in $\mathcal G_m$ uses all four directions. By
Step 1, the skeletons with all four directions that are not in $\mathcal G_m$, with $m\in J_N$,
contribute at most $\sum_m\mathrm{bin}(m)\,25m^2\,C/m^2\le25C$ to $N^2o^{(4)}_N$.

\emph{Step 3 (run lengths given the skeleton).} We show that $N^2\varpi_N=g(a,b,c,d)(1+o(1))$
uniformly over the skeletons in $\mathcal G_m$ with $m\in J_N$. For large $N$ such $m$ satisfy
$m_0\le m\le2\beta N^{1/8}\le N^{1/6}$. On $\mathcal G_m$ we have $A,B\ge m/4$, so
$A-1,B-1\ge m/8$. Put $f(X)=X^{A-2}(\nu-X)^{B-2}$. For
$X\ge1$ we have $\chi_a(X)=\frac{X^{a-1}}{(a-1)!}\prod_{1\le j<a}(1-j/X)_+$, where
$(y)_+=\max(y,0)$, and similarly for $b,c,d$. The terms $X=0$ and $X=\nu$ of $\Omega_N$ vanish. So
\[
 \Omega_N(a,b,c,d)=\frac1{\Gamma_4}\sum_{X=1}^{\nu-1}f(X)\omega(X),
\]
where $\omega(X)\in[0,1]$ is the product of the four products. The function $f$ vanishes at $0$ and
$\nu$ and is unimodal, and its mode $X^*=\nu(A-2)/(m-3)$ lies in $[\nu/6,5\nu/6]$ since
$A,B\ge m/4$ and $m\ge18$. On each monotone piece of $f$ every value $f(X)$ lies between the
integrals of $f$ over the two unit intervals next to $X$. So $\lvert\sum_{j\le X\le k}f(X)-\int_j^kf\rvert\le2\max f$ for all integers
$0\le j\le k\le\nu$. Moreover $f/\int_0^\nu f$ is the density of $\nu$ times a Beta variable with
parameters $A-1$ and $B-1$, and $\int_0^\nu f=\nu^{m-2}\mathrm B(A-1,B-1)$.

We need three estimates. First, this Beta law is the law of the $(A-1)$-th smallest of $m-2$
independent uniform points in $[0,1]$. A union bound over sets of $A-1$ points gives, for $0<t<1$,
\[
 \Prob(\mathrm{Beta}\le t)\le\binom{m-2}{A-1}t^{A-1}\le\Bigl(\frac{e(m-2)t}{A-1}\Bigr)^{A-1}
 \le(8et)^{A-1}.
\]
We take $t=2\nu^{-1/4}$. Since $A-1\ge2$ and $16e\nu^{-1/4}<1$ for large $N$, this gives
$\int_0^{\nu^{3/4}+1}f\le(16e)^2\nu^{-1/2}\int_0^\nu f$, and the same bound holds near $\nu$. Second,
for $\nu^{3/4}\le X\le\nu-\nu^{3/4}$, the inequality $1-\prod_j(1-t_j)_+\le\sum_jt_j$ for $t_j\ge0$
gives $1-\omega(X)\le(a^2+b^2+c^2+d^2)/(2\nu^{3/4})\le m^2\nu^{-3/4}$. Third, on
$[\nu/12,11\nu/12]$ we have
$(\log f)''=-(A-2)/X^2-(B-2)/(\nu-X)^2\ge-144m/\nu^2$. This interval contains every $X$ with
$\lvert X-X^*\rvert\le\nu/(12\sqrt m)$, and for these $X$ Taylor's formula at $X^*$ gives
$f(X)\ge e^{-1/2}\max f$. So $\int_0^\nu f\ge e^{-1/2}\max f\cdot\nu/(6\sqrt m)$, that is,
$\max f\le6e^{1/2}\sqrt m\,\nu^{-1}\int_0^\nu f$.

Now we put the three estimates together. For the upper bound,
$\sum_{X=1}^{\nu-1}f\omega\le\sum_Xf\le\int_0^\nu f+2\max f$. For the lower bound we sum only over
$\nu^{3/4}\le X\le\nu-\nu^{3/4}$, use the second estimate for $\omega$, compare the sum of $f$ with
its integral, and use the first estimate for the two ends of the integral. This gives
\[
 \sum_{X=1}^{\nu-1}f\omega\ge\bigl(1-m^2\nu^{-3/4}\bigr)
 \Bigl(\int_0^\nu f-2(16e)^2\nu^{-1/2}\int_0^\nu f-2\max f\Bigr).
\]
Since $m^2\le N^{1/3}$ and $\sqrt m/\nu\le N^{-5/6}$, we get
$\Gamma_4\Omega_N=\nu^{m-2}\mathrm B(A-1,B-1)(1+o(1))$, uniformly. Also
$\binom{N-1}m=\frac{N^m}{m!}\prod_{j=1}^m(1-\frac jN)=\frac{N^m}{m!}(1+O(m^2/N))$. With
$\nu^{m-2}/N^{m-2}=2^{-(m-2)}$ and $\mathrm B(A-1,B-1)=(A-2)!\,(B-2)!/(m-2)!$ this gives
$N^2\varpi_N=g(a,b,c,d)(1+o(1))$.

\emph{Step 4 (average over the skeleton).} We show that
$\E[g\,\mathbf 1_{\mathcal G_m}\mid m]/m\to c_*=\frac{2(r+s)}{\pi\beta}$ as $m\to\infty$. First we
need a local limit for the fair binomial law. Let $0\le k\le n$ and $y=(2k-n)/n$. If
$\lvert2k-n\rvert\le n^{3/5}$, Stirling's formula gives
\begin{align*}
 \sqrt n\binom nk2^{-n}&=\sqrt{\frac2{\pi(1-y^2)}}\,
 e^{-\frac n2[(1+y)\log(1+y)+(1-y)\log(1-y)]}\bigl(1+O(\tfrac1n)\bigr)\\
 &=\sqrt{\frac2\pi}\,e^{-\frac{(2k-n)^2}{2n}}\bigl(1+O(n^{-3/5})\bigr),
\end{align*}
since the bracket is $y^2+O(y^4)$ and $ny^4\le n^{-3/5}$. Both sides decrease as $\lvert2k-n\rvert$
grows, so outside this range both are $O(e^{-n^{1/5}/3})$. Hence
\begin{equation}\label{eq:app-planar-dml}
 \sup_{0\le k\le n}\Bigl\lvert\sqrt n\binom nk2^{-n}-\sqrt{\tfrac2\pi}\,e^{-(2k-n)^2/(2n)}\Bigr\rvert
 =O(n^{-1/2}).
\end{equation}
On $\mathcal G_m$ we have $A-2,B-2\ge m/8$. We apply \eqref{eq:app-planar-dml} with $n=A-2$,
$k=a-1$ and with $n=B-2$, $k=c-1$, and note that $2(a-1)-(A-2)=a-b$. Then, uniformly on
$\mathcal G_m$,
\[
 \frac gm=\frac{m-1}{\pi\sqrt{(A-2)(B-2)}}
 \exp\Bigl(-\frac{(a-b)^2}{2(A-2)}-\frac{(c-d)^2}{2(B-2)}\Bigr)+O(m^{-1/2}).
\]
Now we replace $m-1$ by $m$, then $A-2$ by $A$ and $B-2$ by $B$, and finally $B$ by $m(1-\theta)$,
where $\theta=A/m$. The last step is needed because $B=m(1-\theta)+1$. Each step changes the
prefactor by a factor $1+O(1/m)$ and each exponential by $O(1/m)$, since
$0\le e^{-y/A_2}-e^{-y/A_1}\le\frac{A_2-A_1}{eA_1}$ for $y\ge0$ and $0<A_1\le A_2$. So, uniformly on
$\mathcal G_m$,
\[
 \frac gm=\Psi\Bigl(\frac{a-b}{\sqrt m},\frac{c-d}{\sqrt m},\frac Am\Bigr)+O(m^{-1/2}),\qquad
 \Psi(x_1,x_2,\theta)=\frac{e^{-x_1^2/(2\theta)-x_2^2/(2(1-\theta))}}{\pi\sqrt{\theta(1-\theta)}}.
\]
On $\mathcal G_m$ we have $\theta\in[\frac14,\frac34+\frac1m]\subset[\frac15,\frac45]$, as $m\ge20$.
Let $\tilde\Psi$ be $\Psi$ with $\theta$ replaced by the nearest point of $[\frac15,\frac45]$. It is
bounded and continuous on $\R^3$ and equals $\Psi$ on $\mathcal G_m$.

Next we prove a central limit theorem for $W_m=(a-b)+\mathrm i(c-d)$. Put
$\zeta_j=\mathrm i^{\xi_j}+\varkappa$, where $\xi_j=Z_j-Z_{j-1}$. The $\zeta_j$ are independent with
$\E\zeta_j=0$, $\E\lvert\zeta_j\rvert^2=1-\varkappa^2$ and
$\E\zeta_j^2=\vartheta_2-\varkappa^2=-(1-\varkappa)^2$. Since
$\mathrm i^{Z_j}=\mathrm i^{Z_{j-1}}(\zeta_j-\varkappa)$,
\[
 (1+\varkappa)W_m=M_m+\mathrm i^{Z_0}+\varkappa\,\mathrm i^{Z_m},\qquad
 M_m=\sum_{j=1}^m\mathrm i^{Z_{j-1}}\zeta_j,
\]
and $M_m$ is a martingale with increments of modulus at most $2$. We take the skeletons for all $m$
to be the first steps of one infinite walk, so the $\sigma$-fields are nested. Fix $\varphi\in\R$. For a
constant $\omega$ with $\lvert\omega\rvert=1$ we have
$\E(\operatorname{Re}\omega\zeta_j)^2=\frac12(\E\lvert\zeta_j\rvert^2+\operatorname{Re}\omega^2\E\zeta_j^2)$.
So the $j$-th increment of $\operatorname{Re}(e^{-\mathrm i\varphi}M_m)$ has conditional variance
$\frac12[(1-\varkappa^2)-(1-\varkappa)^2\cos(2\varphi)(-1)^{Z_{j-1}}]$. The average of these
variances tends to $\frac12(1-\varkappa^2)$ in probability, because
$\sum_{t=0}^m(-1)^{Z_t}=V_2$ and $\E\lvert V_2/m\rvert^4\le C/m^2$ by Step 2. The increments are
bounded, so the martingale central limit theorem \cite[Cor.~3.1]{hallheyde1980} and the
Cram\'er--Wold device give $M_m/\sqrt m\Rightarrow N(0,\frac12(1-\varkappa^2)I_2)$. Hence
\[
 \frac{W_m}{\sqrt m}\Rightarrow X\sim N(0,\sigma_*^2I_2),\qquad
 \sigma_*^2=\frac{1-\varkappa}{2(1+\varkappa)}=\frac r{2(r+s)}.
\]
Also $A=\frac{m+1}2+\frac12V_2$, so $A/m\to\frac12$ in probability, and
$(W_m/\sqrt m,A/m)\Rightarrow(X,\frac12)$. By Step 2,
$\E[\tilde\Psi\,\mathbf 1_{\mathcal G_m^c}\mid m]\le C/m^2$. So, since $\tilde\Psi$ is bounded and
continuous,
\[
 \frac{\E[g\,\mathbf 1_{\mathcal G_m}\mid m]}m
 =\E\,\tilde\Psi\Bigl(\frac{W_m}{\sqrt m},\frac Am\Bigr)+O(m^{-1/2})
 \to\E\,\Psi\bigl(X,\tfrac12\bigr)=\frac2\pi\,\E e^{-\lvert X\rvert^2}
 =\frac2{\pi(1+2\sigma_*^2)}=c_*.
\]

\emph{Step 5 (conclusion).} By Steps 1 and 2,
$N^2o^{(4)}_N=\sum_{m\in J_N}\mathrm{bin}(m)\,\E[N^2\varpi_N\mathbf 1_{\mathcal G_m}\mid m]+O(1)$.
By Step 3 the expectation is $\E[g\,\mathbf 1_{\mathcal G_m}\mid m](1+o(1))$, with $o(1)$ uniform in
$m\in J_N$. Since $\min J_N\to\infty$, Step 4 turns this into
$N^2o^{(4)}_N=c_*(1+o(1))\sum_{m\in J_N}\mathrm{bin}(m)\,m+O(1)$. As in Step 1,
$\E[m\mathbf 1\{m\notin J_N\}]\to0$, and $\bar m=\beta T(1+O(1/N))$. Hence
$N^2o^{(4)}_N=c_*\beta T(1+o(1))+O(1)=\frac{2(r+s)}\pi T(1+o(1))$.
\end{proof}

\subsection{Turns only}

We prove parts (b) and (c) of Theorem~\ref{thm:pureturns}. Throughout, a word is a word of length
$n$ in the letters $\pm$, its weight is $x^j$ if it has $j$ switches, and we say that it switches at
step $t$ if its letters $t$ and $t+1$ differ ($1\le t\le n-1$). The words of length $n$ have total
weight $2(1+x)^{n-1}$, and those with $k$ letters $+$ have total weight $S_{k,n-k}(x)$ if
$1\le k\le n-1$ and $1$ if $k\in\{0,n\}$. We need one fact about Bessel functions. Let
$F(y)=e^{-y}(I_0(y)+I_1(y))$ for $y>0$, so that $\psi(y)=ye^yF(y)$. The integral forms
$I_n(y)=\frac1\pi\int_0^\pi e^{y\cos\theta}\cos(n\theta)\,d\theta$ of $I_0$ and
$I_1$~\cite[Eq.~10.32.3]{dlmf} and the substitution $\varsigma=1-\cos\theta$ give
\[
 F(y)=\frac1\pi\int_0^\pi e^{-y(1-\cos\theta)}(1+\cos\theta)\,d\theta
 =\frac1\pi\int_0^2e^{-y\varsigma}\sqrt{\frac{2-\varsigma}\varsigma}\,d\varsigma .
\]
So $F$ is positive and decreasing. Replacing $y$ by $y/2$ and $\varsigma$ by $2\varsigma$ gives
\[
 F(y/2)=\frac2\pi\int_0^1e^{-y\varsigma}\sqrt{\frac{1-\varsigma}\varsigma}\,d\varsigma .
\]
This is at most $2F(y)$, because $1-\varsigma\le2-\varsigma$ and the part of the integral for
$F(y)$ over $[1,2]$ is nonnegative. Hence
\begin{equation}\label{eq:app-planar-F}
 F(y')\le2F(y)\qquad\text{whenever }y/2\le y'\le y .
\end{equation}

\begin{lemma}[one switch of a balanced word]\label{lem:oneswitch}
Let $N\ge4$ be even, let $x>0$, and put $z=Nx$ and $\eta=z(z+2)/(2N)$. For $1\le t\le N-1$ let
$S^{(t)}_N(x)$ be the total weight of the balanced words of length $N$ that switch at step $t$. If
$\eta\le\frac12$, then $S^{(t)}_N(x)\le2x+8x\,S_N(x)$ for every $t$, and
$\sum_tS^{(t)}_N(x)=xS_N'(x)\le(2z+2)S_N(x)$.
\end{lemma}

In words, a balanced word switches at a given step with probability of order $x$, as a word
without any condition does, and it has at most about $2z$ switches on average.

\begin{proof}
Reversing a word maps the balanced words that switch at $t$ to those that switch at $N-t$, and it
keeps the weight. So we may take $t\ge N/2$. A balanced word that switches at $t$ is a word $w_1$
of length $t$ followed by a word $w_2$ of length $N-t$ whose first letter differs from the last
letter of $w_1$, and its weight is $x$ times the product of the two weights. If $w_1$ has $k$
letters $+$, then $w_2$ has $N/2-k$ of them, so $k\le N/2\le t$ and $t-k\le N/2\le t$.

Consider first the balanced words with $1\le k\le t-1$. Fix $w_2$. Then $k$ is fixed, and the
words $w_1$ that can precede $w_2$ have total weight at most $S_{k,t-k}(x)$. By
\cite[Proposition~4.1]{paperA} the bin polynomials grow coefficientwise as the bin moves toward the
middle, that is, $S_{a,b}\le S_{a+1,b-1}$ coefficientwise when $a<b-1$. Since also
$S_{a,b}=S_{b,a}$, this gives $S_{k,t-k}(x)\le S_t(x)$. The words $w_2$ of length $N-t$ have total
weight at most $2(1+x)^{N-t-1}\le2e^{(N-t)x}$. By~\eqref{eq:planar-bessel} with $t\ge2$ steps,
$S_t(x)\le\frac2t\psi(tx)=2xe^{tx}F(tx)$. So these words contribute at most
\[
 x\cdot2xe^{tx}F(tx)\cdot2e^{(N-t)x}=4x^2e^{z}F(tx)\le8x^2e^zF(z)=4x\cdot\frac2N\psi(z)
 \le\frac{4x}{1-\eta}\,S_N(x)\le8x\,S_N(x),
\]
to $S^{(t)}_N(x)$. Here the first inequality is~\eqref{eq:app-planar-F} with $z/2\le tx\le z$, and
the next one is the lower bound in~\eqref{eq:planar-bessel}. If $k=0$ or $k=t$, then $w_1$ is
constant and $t=N/2$, and $w_2$ is the constant word of the other letter. These two words have
weight $x$ each. This proves the first claim.

Each balanced word with $j$ switches is counted $j$ times in $\sum_tS^{(t)}_N$, so the sum is
$xS_N'(x)$. The function $y\mapsto\log S_N(e^y)$ is convex, since $S_N$ has nonnegative
coefficients. So its derivative at $y=\log x$ is at most its increase over $[\log x,\log x+1]$,
that is, $xS_N'(x)/S_N(x)\le\log(S_N(ex)/S_N(x))$. By~\eqref{eq:planar-bessel},
$S_N(ex)\le\frac2N\psi(ez)$ and $S_N(x)\ge\frac12\cdot\frac2N\psi(z)$, and
$\psi(ez)/\psi(z)=e\cdot e^{(e-1)z}F(ez)/F(z)\le e^{1+(e-1)z}$. Hence
$xS_N'(x)/S_N(x)\le1+\log2+(e-1)z\le2z+2$.
\end{proof}

\begin{proof}[Proof of Theorem~\ref{thm:pureturns}(b)]
A pair of balanced words with a common switch has a common switch at some step $t$, and the pairs
that both switch at $t$ have total weight $S^{(t)}_N(x)^2$. So
$\Delta_N(x)\le\sum_tS^{(t)}_N(x)^2\le\max_tS^{(t)}_N(x)\cdot\sum_tS^{(t)}_N(x)$, and
Lemma~\ref{lem:oneswitch} gives the bound, because $z(z+2)\le N$ means $\eta\le\frac12$.
\end{proof}

\begin{proof}[Proof of Theorem~\ref{thm:pureturns}(c)]
We compare $v^*_N$ with a value $x_1$ of the odds slightly above $u_N$, at which the bound of
part (b) is already smaller than the gain in $S_N^2$.

Put $z_0=Nu_N$. Since $S_N(u)\ge2u$, we have $u_N\le\frac12$ and $p_N\ge\frac23$. With
$z_0=z_N/p_N$ and $z_N\le L$~\cite[Proposition~3.8]{paperA} this gives $z_0\le\frac32L$. By
\eqref{eq:planar-bessel} and $S_N(u_N)=1$, $\psi(z_0)\ge N/2$. Let $\kappa=40(z_0+2)^2/N$, which
tends to $0$, and let $x_1>u_N$ be the root of $S_N(x_1)=1+\kappa$, with $z_1=Nx_1$ and
$\eta_1=z_1(z_1+2)/(2N)$.

We first show that $z_1\le z_0+1$ for large $N$. Since $(\log\psi)'\ge1$, we have
$\psi(z_0+1)\ge e\,\psi(z_0)\ge eN/2$. So by~\eqref{eq:planar-bessel}, at $x=(z_0+1)/N$ we get
$S_N(x)\ge(1-O(L^2/N))\,e$, which exceeds $1+\kappa$ for large $N$. As $S_N$ is increasing,
$z_1\le z_0+1$. In particular $z_1(z_1+2)\le N$ and $\kappa\le1$ for large $N$, so part (b) gives
\[
 \Delta_N(x_1)\le\bigl(2x_1+16x_1\bigr)(2z_1+2)\cdot2=\frac{72\,z_1(z_1+1)}N
 \le\frac{72(z_0+1)(z_0+2)}N\le2\kappa .
\]
By part (a), $o_N/c_N=(1+\kappa)^2-\Delta_N(x_1)\ge1$ at $v=x_1$. The ratio $o_N/c_N$ increases
with $v$, so $v^*_N\le x_1$. Now $\psi(z_1)\le\frac N2(1+\kappa)/(1-\eta_1)$
by~\eqref{eq:planar-bessel}, and $(\log\psi)'\ge1$. Hence
\[
 0\le Nv^*_N-z_0\le z_1-z_0\le\log\psi(z_1)-\log\psi(z_0)\le\log\frac{1+\kappa}{1-\eta_1}
 =O\Bigl(\frac{L^2}N\Bigr),
\]
where the first inequality is part (a). This is the first claim. For the second,
$v=rh/(1-2rh)$ gives $rT^*_N=Nv^*_N/(1+2v^*_N)$, and $z_N=z_0/(1+u_N)$. Both $v^*_N$ and $u_N$ are
$O(L/N)$, so $rT^*_N=Nv^*_N+O(L^2/N)$ and $z_N=z_0+O(L^2/N)$.
\end{proof}
